\section{Câu 4}
Cho mạch khuếch đại tín hiệu như hình vẽ. Các tụ $C_1$, $C_2$ có giá trị rất lớn. BJT có hệ số $\beta = 100$ và có mã là 2N2907

\begin{figure}[h!]
    \centering
    \resizebox{0.7\textwidth}{!}{
    \begin{tikzpicture}[transform shape]
	% Paths, nodes and wires:
	\draw (9.25, 5.25) to[american resistor, mirror, l_={$R_3=10k \Omega$}] (9.25, 6.75);
	\draw (8, 6.25) to[american resistor, mirror, l_={$R_1 = 20k\Omega$}] (8, 4.75);
	\node[npn] at (12, 8){};
	\node[ocirc](N1) at (14, 4.25){} node[anchor=west] at (N1.east){$v_o$};
	\node[ground] at (8, 1.75){};
	\draw (5.15, 8) to[american resistor, mirror, l_={$R_{S}=10k \Omega$}] (3.75, 8);
	\node[oscillator](N2) at (2.49, 6.49){} node[anchor=west] at (N2.east){$v_s$};
	\draw (8, 3.75) to[american resistor, mirror, l_={$R_2 = 20k\Omega$}] (8, 2.25);
	\node[ground] at (12, 1.75){};
	\node[ground] at (2, 5){};
	\draw (6, 8) to[capacitor, l={$\infty$}] (7, 8);
	\node[vcc](N3) at (8, 6.5){} node[anchor=south] at (N3.text){$+9V$};
	\draw (2, 7) -| (2, 8) -- (3.75, 8);
	\draw (5.15, 8) -- (6, 8);
	\draw (7, 8) -- (11.16, 8);
	\draw (8, 6.5) -| (8, 6.25);
	\draw (9.25, 6.75) -- (9.25, 8);
	\draw (8, 4.75) -| (8, 3.75);
	\draw (8, 4.25) -- (9.25, 4.25) -- (9.25, 5.25);
	\draw (10, 4.25) to[capacitor, l={$C_B = \infty$}] (11, 4.25);
	\draw (12, 4) to[american resistor, mirror, l_={$R_4 = 2k\Omega$}] (12, 2.5);
	\draw (9.25, 4.25) -- (10, 4.25);
	\draw (11, 4.25) -| (12, 4);
	\draw (12, 7.25) -| (12, 4.25);
	\node[vcc](N4) at (12, 9.5){} node[anchor=south] at (N4.text){$+9V$};
	\draw (12, 9.5) -| (12, 8.77);
	\draw (12, 4.25) -- (14, 4.25);
	\draw (12, 2.5) -| (12, 1.75);
	\draw (8, 2.25) -| (8, 1.75);
	\draw (2, 6) -| (2, 5);
	\node[circ] at (9.25, 8){};
	\node[circ] at (8, 4.25){};
	\node[circ] at (9.25, 4.25){};
	\node[circ] at (12, 4.25){};
\end{tikzpicture}
    }
\end{figure}
\textbf{a) Vẽ VTC của mạch (kiểm chứng sử dụng mô phỏng) và tìm điểm hoạt động $Q$ của BJT.}

Theo DC mạch có dạng:
\begin{figure}[htbp]
\centering
\begin{subfigure}{0.48\textwidth}
  \centering
  \begin{tikzpicture}[transform shape]
	% Paths, nodes and wires:
	\draw (8.75, 8) to[american resistor, mirror, l_={$R_3=10k \Omega$}] (10.25, 8);
	\draw (7.5, 10) to[american resistor, mirror, l_={$R_1 = 20k\Omega$}] (7.5, 8.5);
	\node[npn] at (12, 8){};
	\node[ocirc](N1) at (14, 6.75){} node[anchor=west] at (N1.east){$v_o$};
	\node[ground] at (7.5, 4.5){};
	\draw (7.5, 7.5) to[american resistor, mirror, l_={$R_2 = 20k\Omega$}] (7.5, 6);
	\node[ground] at (12, 4.5){};
	\node[vcc](N2) at (7.5, 10.25){} node[anchor=south] at (N2.text){$+9V$};
	\draw (7.5, 10.25) -| (7.5, 10);
	\draw (7.5, 8.5) -| (7.5, 7.5);
	\draw (12, 6.5) to[american resistor, mirror, l_={$R_4 = 2k\Omega$}] (12, 5);
	\node[vcc](N3) at (12, 9.5){} node[anchor=south] at (N3.text){$+9V$};
	\draw (12, 9.5) -| (12, 8.77);
	\draw (12, 6.75) -- (14, 6.75);
	\draw (12, 5) -| (12, 4.25);
	\draw (7.5, 6) -- (7.5, 4.5);
	\node[circ] at (7.5, 8){};
	\node[circ] at (12, 6.75){};
	\draw (7.5, 8) -- (8.75, 8);
	\draw (10.25, 8) -- (11.16, 8);
	\draw (12, 7.23) -| (12, 6.5);
\end{tikzpicture}
%   \includegraphics[width=\linewidth]{fig1.png}
    \caption{Mạch ở chế dộ DC}
%   \label{fig:first}
\end{subfigure}
\hfill
\begin{subfigure}{0.48\textwidth}
  \centering
  \begin{tikzpicture}[transform shape]
	% Paths, nodes and wires:
	\draw (8.75, 8) to[american resistor, mirror, l={$R_{th}$}] (10.25, 8);
	\node[npn] at (12, 8){};
	\node[ocirc](N1) at (14, 6.75){} node[anchor=west] at (N1.east){$v_o$};
	\node[ground] at (12, 4.5){};
	\node[vcc](N2) at (8, 9.5){} node[anchor=south] at (N2.text){$V_{th}$};
	\draw (8, 9.5) -| (8, 9.25);
	\draw (12, 6.5) to[american resistor, mirror, l_={$R_4 = 2k\Omega$}] (12, 5);
	\node[vcc](N3) at (12, 9.5){} node[anchor=south] at (N3.text){$+9V$};
	\draw (12, 9.5) -| (12, 8.77);
	\draw (12, 6.75) -- (14, 6.75);
	\draw (12, 5) -| (12, 4.25);
	\node[circ] at (12, 6.75){};
	\draw (10.25, 8) -- (11.16, 8);
	\draw (12, 7.23) -| (12, 6.5);
	\draw (8, 9.25) -| (8, 8) -- (8.75, 8);
\end{tikzpicture}
%   \includegraphics[width=\linewidth]{fig2.png}
    \caption{Mạch sau khi đã Thevenin}
%   \label{fig:second}
\end{subfigure}
% \caption{Overall caption for both}
% \label{fig:both}
\end{figure}

Thevenin ta có:
\[ R_{th} = R_{1} // R_{2} + R_3 = 20\,\textsf{k}\Omega \]
\[ V_{th} = 9\times\dfrac{R_{2}}{R_{1} + R_{2}}= 4.5 \,\textsf{V} \]

Áp dụng KVL cho vòng từ cực B tới cực E:
\[ I_{B}R_{th} + V_{BE} + I_{E}R_{4} - 9 = 0 \]

Ta có: $I_E = (\beta + 1)I_B$
\[
\Longrightarrow I_B = \frac{9 - V_{BE}}{R_{th} + (\beta + 1)R_4}
= \frac{9 - 0.7}{20k + (100 + 1)\times 2k}
= 0.0171\,\text{mA}
\]

\[
I_C = \beta I_B = 100 \times 0.0171 = 1.7117\,\text{mA}
\]

Áp dụng KVL cho vòng từ cực C tới cực E:
\[-9+V_{CE}+I_E R_4-9=0 \]
\[\Rightarrow V_{CE}=18-(\beta+1)I_B\times R_4=5.54\ V\]

Vậy điểm làm việc $Q$ là: $Q(I_{CQ},V_{ECQ}) = Q(1.7117\ mA,\ 5.54\ V)$

\textbf{b) Đặt $v_{s} = V_{m} \sin \left( \omega t\right)$ vào mạch. Tìm $A_{vo}$, $G_v$, $R_i$, $R_o$ của mạch.}

Mạch ở chế độ AC có dạng:
\begin{figure}[h!]
    \centering
    \resizebox{0.7\textwidth}{!}{
\begin{tikzpicture}[transform shape]
	% Paths, nodes and wires:
	\draw (9.25, 5.75) to[american resistor, mirror, l_={$R_3=10k \Omega$}] (9.25, 7.25);
	\node[npn] at (12, 8){};
	\node[ocirc](N1) at (16, 5){} node[anchor=west] at (N1.east){$v_o$};
	\draw (7.15, 8) to[american resistor, mirror, l_={$R_{S}=10k \Omega$}] (5.75, 8);
	\node[oscillator](N2) at (4.49, 6.49){} node[anchor=west] at (N2.east){$v_s$};
	\node[ground] at (12, 3){};
	\node[ground] at (4, 5){};
	\draw (4, 7) -| (4, 8) -- (5.75, 8);
	\draw (7.15, 8) -- (8, 8);
	\draw (12, 4.75) to[american resistor, mirror, l_={$R_E$}] (12, 3.25);
	\draw (12, 5) -- (16, 5);
	\draw (4, 6) -| (4, 5);
	\node[circ] at (9.25, 8){};
	\node[circ] at (12, 5){};
	\node[ground] at (14, 8.75){};
	\draw (12, 8.77) -| (12, 9) -- (14, 9) -| (14, 8.75);
	\draw (8, 8) -| (11.16, 8);
	\draw (9.25, 8) -| (9.25, 7.25);
	\draw (9.25, 5.75) -| (9.25, 5) -- (12, 5) -- (12, 7.25);
	\draw (12, 4.75) -| (12, 5);
	\draw (12, 3) -| (12, 3.25);
	\node[shape=rectangle, minimum width=0.715cm, minimum height=1.715cm](N3) at (12.375, 4){} node[anchor=west] at (N3.text){$R_1//R_2//R_4$};
	\draw (15.25, 3.25) to[american resistor, mirror, l_={$R_L=1k \Omega$}] (15.25, 4.75);
	\draw (15.25, 5) -| (15.25, 4.75);
	\node[ground] at (15.25, 3){};
	\draw (15.25, 3.25) -| (15.25, 3);
\end{tikzpicture}
}
\end{figure}

Xét hoạt động chế độ AC cho toàn mạch:
\begin{equation*}
	R_E = R_1//R_2//R_4 \approx 1.6667\ (k\Omega)
\end{equation*}
\begin{equation*}
	g_m = \dfrac{I_C}{V_T} = \dfrac{1.7117m}{25m} = 64.468\ (mS)
\end{equation*}
\[
r_\pi = \frac{V_T}{I_B} = \frac{\beta V_T}{I_C} = 1461\ \Omega
\]

\[
r_e = \frac{V_T}{I_E} = \frac{V_T}{(\beta+1)I_B} = 14.475\ \Omega
\]

\[
R_{in} = \left. \frac{v_i}{i_i} \right|_{i_o=0} = \dfrac{v_{be} + v_o}{i_{R_3}+i_e-\alpha i_e}
\]

Trong đó:
\begin{itemize}
	\item[+] $i_{R_3} = \dfrac{v_{R_3}}{R_3}=\dfrac{v_{be}}{R_3}=\dfrac{i_e r_e}{R_3}$
	\item[+] $v_i = v_{be} + v_o = i_e r_e + R_E(i_{R_3}+i_e)=i_e r_e + R_E\left(\dfrac{i_e r_e}{R_3}+i_e\right)$
\end{itemize}

\begin{equation*}
	\Rightarrow R_{in} = \dfrac{i_e r_e + R_E\left(\dfrac{i_e r_e}{R_3}+i_e\right)}{\dfrac{i_e r_e}{R_3}+i_e-\alpha i_e} = \dfrac{r_e +  R_E\left(\dfrac{r_e}{R_3}+1\right)}{\dfrac{r_e}{R_3}+1-\alpha} = 
\end{equation*}

\[
A_{vo} = \left. \frac{v_o}{v_i} \right|_{R_L=\infty} = \dfrac{R_E\left(\dfrac{i_e r_e}{R_3}+i_e\right)}{i_e r_e + R_E\left(\dfrac{i_e r_e}{R_3}+i_e\right)} = \dfrac{R_E\left(\dfrac{r_e}{R_3}+1\right)}{r_e + R_E\left(\dfrac{r_e}{R_3}+1\right)}
\]

\[
R_o = \left. \frac{v_o}{i_o} \right|_{v_i=0}= \dfrac{v_o}{i_{R_3}+i_e-i_{R_E}} = \dfrac{R_3}{\beta+1}//r_e//R_E
\]

\[
G_v = \frac{R_{in}}{R_{in} + R_{sig}} \times 
\frac{R_L}{R_L + R_o} \times A_{vo} 
= \frac{4.52}{4.52 + 1} \times 
\frac{300}{300 + 1.65} \times 0.997 
= 0.81\ V/V
\]



\textbf{c) Lựa chọn các tụ $C_{1}$, $C_{2}$ để mạch có $f_{L}=100\,\textsf{Hz}$.}